\documentclass[12pt,letterpaper]{article}
\usepackage{amsmath}
\usepackage{amsthm}
\usepackage{amsfonts}
\usepackage{amssymb}
\usepackage{amscd}
\usepackage{enumerate}
\usepackage{fancyhdr}
\usepackage{mathrsfs}
\usepackage{bbm}
\usepackage{framed}
\usepackage{mdframed}
\usepackage{cancel}
\usepackage{mathtools}
\usepackage{verbatim}
\usepackage{hyperref}
\usepackage[letterpaper,voffset=-.5in,bmargin=3cm,footskip=1cm]{geometry}
\setlength{\parindent}{0.0in}
\setlength{\parskip}{0.1in}
\allowdisplaybreaks
\headheight 15pt
\headsep 10pt
\newcommand\N{\mathbb N}
\newcommand\Z{\mathbb Z}
\newcommand\R{\mathbb R}
\newcommand\Q{\mathbb Q}
\newcommand\lcm{\operatorname{lcm}}
\newcommand\setbuilder[2]{\ensuremath{\left\{#1\;\middle|\;#2\right\}}}
\newcommand\E{\operatorname{E}}
\newcommand\V{\operatorname{V}}
\newcommand\Pow{\ensuremath{\operatorname{\mathcal{P}}}}

\DeclarePairedDelimiter\ceil{\lceil}{\rceil}
\DeclarePairedDelimiter\floor{\lfloor}{\rfloor}
\newcommand\hint[1]{\textbf{Hint}: #1}
\newcommand\note[1]{\textbf{Note}: #1}
\newcounter{enum22i}
\newenvironment{22enumerate}{\begin{enumerate}[a.]\itemsep0em\setcounter{enumi}{\value{enum22i}}}{\setcounter{enum22i}{\value{enumi}}\end{enumerate}}
\newenvironment{22itemize}{\begin{itemize}\itemsep0em}{\end{itemize}}
\fancypagestyle{firstpagestyle} {
  \renewcommand{\headrulewidth}{0pt}
  \lhead{\textbf{CSCI 0220}}
  \chead{\textbf{Discrete Structures and Probability}}
  \rhead{M. Littman}
}

\fancypagestyle{fancyplain} {
  \renewcommand{\headrulewidth}{0pt}
  \lhead{\textbf{CSCI 0220}}
  \chead{Homework 2}
}
\pagestyle{fancyplain}
\usepackage{graphicx}
\usepackage{graphicx,amsmath}

\begin{document}
  \thispagestyle{firstpagestyle}
  \begin{center}
  \vspace*{7px}
    {\huge \textbf{Homework 2}}

    {\large Due: \textit{Monday, May 31$^{st}$ at 11:59pm EDT}}
  \end{center}

    All homeworks are due the Monday after their release at 11:59pm EDT on Gradescope.

    Please do not include any identifying information about yourself in the handin, including your Banner ID.

    Be sure to fully explain your reasoning and show all work for full credit. We recommend checking out the sample proofs on the course website to see the level of rigor for which we're looking.

\subsection*{Problem 1}
{
\nopagebreak 

{\setcounter{enum22i}{0}
  
	Prove each of the following by contrapositive:
	\begin{22enumerate}
		\item If $x^{3}$ is even, then $x$ is even.
		\item If $a + b \geq 25$, where $a$ and $b$ are both
			integers, then either $a \geq 13$ or $b \geq 13$.
	\end{22enumerate}
	Prove the following claim via proof by contradiction: 
	\begin{22enumerate}
    	 \item For all real numbers $x$ and $y$, if $x$ is a non-zero rational number and $y$ is an irrational number, then $\frac{x}{y}$ is irrational.
	\end{22enumerate}
}
}


\subsection*{Problem 2}
{
\nopagebreak 

{\setcounter{enum22i}{0}

      \begin{22enumerate}

      \item
      Having proudly won planet status centuries ago, our smallest planet Mercury has spent many orbits trying to determine which of its fellow near-Sun orbiters could also be considered planets. To figure this out, it decides to create a poll to determine what its nearest planet friends, Venus, the Earth and Mars, think. It plans on starting the vote off with \underline{\href{https://en.wikipedia.org/wiki/2005_HC4}{2005 HC4}}. However, it needs a machine that can handle the data.

      To help it out, draw a circuit using only two-input \textsc{AND}, \textsc{OR}, and \textsc{NOT} gates that takes the three planets' answers (1 if in favor of HC4 being a planet, 0 if against), and returns the answer of the majority. Be sure to explain how your circuit works.
      
      \item
      Although we have been using \textsc{AND} and \textsc{OR} gates with only two inputs, it is possible to create \textsc{AND} and \textsc{OR} gates that take in more than two inputs. Demonstrate this fact by drawing a circuit that represents a three-input \textsc{AND} and a circuit that represents a three-input \textsc{OR} gate using only two-input \textsc{AND} and \textsc{OR} gates.

      \item
      Thanks to your help, Mercury has started running the polls. However, as Mercury considers the magic behind your circuit, it realizes how little it remembers about numbers. It needs your help to figure out whether the number of votes in favor of any orbiter (0, 1, 2, or 3) is even or odd.

      Draw Mercury a circuit using only \textsc{AND}, \textsc{OR}, and \textsc{NOT} gates to determine whether these four numbers are even or odd. The circuit should have three input wires and one output wire. If an even number of inputs are on (have a value of 1), the output should be off (have a value of 0). If an odd number of input wires are on, the output should be on. You may now use \textsc{AND} and \textsc{OR} gates with $n \geq 2$ inputs (drawn with $n$ input wires and one output wire). Be sure to explain how your circuit works.

      \note The output depends only on \underline{how many} of the input wires are on. For example, if exactly one input is on, then the output should be on. It does not matter which input is on.

      \end{22enumerate}
    
}

}

\subsection*{Problem 3}
{
\nopagebreak 

{\setcounter{enum22i}{0}


      Mercury has continued playing around with logic gates. It considers circuits with $n \geq 0$ input wires and one output wire. Let the wires of the circuit be represented by the set $I_n = \{w_1, w_2, \dots, w_n\}$, where $w_i$ represents the $i$th wire. Consider some subset $t \subseteq I_n$. We say that $t$ \textit{satisfies} the circuit if the circuit evaluates to true when all input wires in $t$ are set to true and all input wires not in $t$ are set to false.

      Let $T_n$ be the set of all $t \subseteq I_n$ such that $t$ satisfies the circuit. Prove that for any $n \ge 2$, there exists a circuit Mercury can build such that:

      \begin{22enumerate}
      \item $|T_n| = 1$
      \item $|T_n| = 2^{n}-1$
      \item $|T_n| = 2^{n}$
      \item $|T_n| = 2^{n-1}$
      \item $|T_n| = n$
      \end{22enumerate}

      \note{You must use every wire in each circuit. You do not need to draw a diagram of your circuits for this problem. We recommend expressing them as their representative logical expressions.}

}
}

\subsection*{Mind Bender (Extra Credit)}
{
Mercury has started thinking about prime numbers. Finding so many unsolved conjectures, it starts out on a quest to create a circuit that outputs whether a number is prime or not. It starts small, considering primes between 0 and 15. In that range, the primes are 2, 3, 5, 7, 11, and 13.

Recall that $n$-bit numbers can represent a binary number, where the position relative to the rightmost digit represents a power of 2. Hence, the binary number $11$ is equivalent to $2^1 + 2^0=3$ in our decimal system, and $101$ is equivalent to $2^2+2^0=5$. Here is the table of binary numbers in our range of interest to get you comfortable with the concept.
\begin{center}
    \begin{tabular}{c|c}
	Binary & Decimal\\
	\hline
	0000 & 0\\
	0001 & 1 \\
	0010 & 2 \\
	0011 & 3 \\
	0100 & 4\\
	0101 & 5 \\
	0110 & 6 \\
	0111 & 7 \\
	1000 & 8\\
	1001 & 9 \\
	1010 & 10 \\
	1011 & 11 \\
	1100 & 12 \\
	1101 & 13 \\
	1110 & 14 \\
	1111 & 15 \\
	\end{tabular}
\end{center}

Draw a circuit that takes in four input wires $A$, $B$, $C$ and $D$ that outputs whether the binary number $DCBA$ is prime by outputting 1 if it is and 0 if it is not. The circuit may use NOT gates, and AND and OR gates with more than two inputs. Here's the mind-bending part: the number of AND gates you use \textbf{cannot exceed four}. Make sure to explain your reasoning.
}

\end{document}